\documentclass[../main.tex]{subfiles}
\begin{document}
%%%%%%%%%%%%%%%%%%%%%%%%%%%%%%%%%%%%%%%%%%%%%%%%%%%%%%%%%%%%%%%%%%%%%%%%%%%%%%%%%%%%%%%%%%%%%%%%%%%
\chapter{Menus and Keyboard Shortcuts}

Now, you could spend the next couple hours, kicked back in your comfy chair cruising through all
of the menus in OmniGraffle to see what's new or where your favorite command went to. Or, you
could just read through this section and see everything in one convenient---and easy to
search---place.

\section{The OmniGraffle Menu}

The OmniGraffle menu, otherwise known as the \textit{app menu}, is where you go to find out more
about OmniGraffle, set its preferences, enter a license (or find out details about your Omni
Account), and more.

\begin{description}[style=nextline]
  \item[{About OmniGraffle}] Find out about the particular version of OmniGraffle you are using.
  \item[{Check for Updates (\textit{for purchases from the Omni Store only})}] Normally,
    OmniGraffle checks for updated versions of itself regularly based on the Update preferences.
    Use this command to check once right now. If there is an updated version of the software,
    you are prompted to download it.
  \item[{Preferences (\textbf{Command-,})}] Open OmniGraffle's preferences window; you can use
    it to customize the application's behavior to your liking.
  \item[{Keyboard Shortcuts}] Opens a window that lets you change or reset OmniGraffle's
    keyboard shortcuts. You can also export and share your custom keyboard setup, and import
    sets that other users create and share.
  \item[{In-App Purchase (\textit{for purchases from the Mac App Store only})}] If you
    downloaded OmniGraffle 7 from the Mac App Store, you can use this menu to purchase upgrades
    to OmniGraffle Standard or OmniGraffle Pro.
  \item[{Free Trial Mode}] During the 14-day free trial, choose this menu item to switch between
    using the features of OmniGraffle Standard and OmniGraffle Pro. This helps you decide which
    version of OmniGraffle is right for you. After the 14-day trial expires, you can still use
    OmniGraffle to view OmniGraffle files.
  \item[{Licenses (\textit{for purchases from the Omni Store only})}] Show the Licenses dialog,
    where you can review your OmniGraffle software licenses.
  \item[{Unlock OmniGraffle}] Sign in to your Omni Account to associate this copy of OmniGraffle
    with it.
  \item[{Manage Omni Account and Purchases}] Open the Omni Account window, where you can manage
    your Omni Account and your OmniGraffle licenses. This command appears after you have signed
    in for the first time.
  \item[{Services}] The Services menu is a collection of system-wide services offered by other
    apps on your Mac. OmniGraffle plays rather well with macOS, so you should be able to use
    your standard services from within the application.
  \item[{Hide OmniGraffle (\textbf{Command-H})}] Leave OmniGraffle open, but temporarily put
    away its windows while you work in another application.
  \item[{Hide Others (\textbf{Option-Command-H})}] Hide all windows of other applications you
    have open, so that you can concentrate on OmniGraffle.
  \item[{Show All}] Choose this option if you chose to \textbf{Hide Others} and then want to get
    all of the windows of those other applications back onscreen. If none of the windows of
    ``other'' open windows are hidden, this menu option is grayed out.
  \item[{Quit OmniGraffle (\textbf{Command-Q})}] Choose this if you really must part with
    OmniGraffle for now. If you have unsaved documents open, you are prompted to save or discard
    your changes.
\end{description}

\section{The File Menu}

\begin{description}[style=nextline]
  \item[{New (\textbf{Command-N})}] If you have a default template, this creates a new document
    based on it. If you don't have a default template, this opens the template chooser.
  \item[{New Resource}] Create a new OmniGraffle stencil, template, or diagram style, starting
    from the template of your choice.

    \begin{description}[style=nextline]
      \item[{\textbf{New Stencil}}] This opens a new document that, when saved, becomes a stencil
        you can reuse in OmniGraffle or share with friends.
      \item[{\textbf{New Template}}] This opens a new document that, when saved, becomes a
        template you can reuse in OmniGraffle or share with friends.
    \end{description}
  \item[{Resource Browser (\textbf{Shift-Command-N})}] Open the Resource Browser, from which you
    can choose an OmniGraffle document, stencil, or template file to edit. You can also use the
    Resource Browser to set a default template to be used whenever you create a new document.
  \item[{Open (\textbf{Command-O})}] Browse for a document to open. You can open any OmniGraffle
    document, or a file in one of the formats OmniGraffle knows how to import.
  \item[{Open Recent}] Open a document you've recently worked on. Choose \textbf{Clear Menu} to
    remove the list of recent documents, so no one knows you've been working on \textit{Surprise
    Party Invitation.graffle}.
  \item[{Close (\textbf{Command-W})}] Close the frontmost window you have open. If you have
    unsaved changes, you are prompted to save or discard them.
  \item[{Close All (\textbf{Option-Command-W})}] Close all open OmniGraffle windows.
  \item[{Save (\textbf{Command-S})}] Write all of your hard work in the frontmost document onto
    the disk for safekeeping. If you have already saved the document once, the new version
    replaces the old one. If it's the first time you're saving this document, you are prompted
    to browse to a location on the disk and enter a file name.
  \item[{Save As (\textbf{Option-Shift-Command-S})}] Save the current OmniGraffle file as
    another filename and/or to another location.
  \item[{Duplicate (\textbf{Shift-Command-S})}] Make a duplicate copy of the current document.
    The duplicated file uses the same filename as the file you're working on, but tacks on
    \textit{copy} to the filename so you know which one is the clone.
  \item[{Rename}] Change the filename of your current project.
  \item[{Move To}] Move the file to another location on your Mac's hard drive.
  \item[{Move to Trash}] Be careful with this one; it moves the file you are currently working
    on into the Trash. This comes in particularly handy when you're working on multiple
    duplicated files as you tweak and compare styles. With the dupes you don't want, just choose
    \textbf{File} $\blacktriangleright$ \textbf{Move to Trash} and then close the window
    (\textbf{Command-W}).
  \item[{Revert To}] When revisions don't really work out like you had planned, choose this to
    go back to the version of the file that's on the disk; it's kind of like a mass-Undo. If you
    open a file that had some changes in an automatic backup, you can choose this command to go
    back to the last explicitly saved version.
  \item[{Place Image}] Browse for an image file; a shape containing the image is created on the
    canvas at the last place you clicked with the Selection tool.
  \item[{Export (\textbf{Option-Command-E})}] This creates a new file, in a file type of your
    choice, from some portion of the data in your OmniGraffle document. See Exporting Image and
    Graphics Files from OmniGraffle, earlier in this guide for more details.
  \item[{Submit to Stenciltown}] Submit the current file as a stencil to Stenciltown, a place
    for OmniGraffle users to find, share, and download stencils. For more information about
    stencils and Stenciltown, see Using, Curating, and Creating Stencils.
  \item[{Upgrade File}] OmniGraffle 7.16 for Mac and 3.14 for iOS introduce a new file type with
    improved support for storing files with cloud sync services. When you open a file created
    with a previous version of OmniGraffle you'll have the option to upgrade it to the new file
    type using this menu item (an upgrade prompt also appears when you make your first change).

    If you choose not to upgrade, the file stays saved in its existing file type. The option to
    upgrade is still available for later if you change your mind.
  \item[{Change File Type}] Choose whether to save the document as a \textit{single file} (the
    default and recommended file type for greatest compatibility with cloud storage), or a
    \textit{package} (a file type useful for document performance in some niche cases).

    Files created with older versions of OmniGraffle that have not yet been upgraded may display
    an additional \textit{legacy} option in this submenu; changing the file type to one of the
    modern options acts the same as upgrading it.

    To share files saved in a modern format with older versions of OmniGraffle, the legacy file
    types are still available as export options.
\end{description}

\begin{ognote}
  Unless you have a specific need to save a document as a file package, it is recommended to save
  all OmniGraffle files in the new default single file type.
\end{ognote}

\begin{description}[style=nextline]
  \item[{Page Setup (\textbf{Shift-Command-P})}] Show the standard macOS Page Setup sheet. This
    sheet contains several panes of settings for how to represent your document on pages.
  \item[{Print (\textbf{Command-P})}] Show the standard macOS Print sheet. In addition to the
    settings in Page Setup, you have a lot of printing-specific options you can set here.
  \item[{Print this Canvas}] Print the currently selected canvas. If you \textbf{Command-Select}
    multiple canvases, this menu option changes to \textbf{Print Canvases}, allowing you to
    print the selected canvases.
\end{description}

\section{The Edit Menu}

\begin{description}[style=nextline]
  \item[{Undo (\textbf{Command-Z})}] OmniGraffle remembers a list of changes you make to a
    document. Choose this command to move one step backwards in the list, as if you hadn't done
    the last thing in the list. Keep choosing it to keep stepping backwards in time.
  \item[{Redo (\textbf{Shift-Command-Z})}] If the last thing you did was to Undo something, you
    can step forward and do it again, effectively Undoing the Undo.
  \item[{Cut (\textbf{Command-X})}] Remove the current selection and put it on the clipboard, so
    you can paste it somewhere else.
  \item[{Copy (\textbf{Command-C})}] Copy the current selection to the clipboard.
  \item[{Copy As}] Copy the selected Canvas items to the clipboard as one of the following file
    types:

    \begin{description}[style=nextline]
      \item[{PDF}] Put a PDF copy of the current selection on the clipboard for pasting into any
        app that accepts PDF files for viewing, editing, or sending. If the app you paste into
        also supports LinkBack (this includes OmniGraffle itself), you can even go back and edit
        the pasted content in OmniGraffle.
      \item[{PNG}] Copies to the clipboard a PNG image asset of the current selection; you can
        then paste this into any app that accepts PNG image files.
      \item[{TIFF}] Copies to the clipboard a TIFF image asset of the current selection; you can
        then paste this into any app that accepts TIFF image files.
      \item[{SVG}] Copies to the clipboard the SVG data for the selected item; you can then paste
        the SVG data into a text editor and save as an SVG file.
      \item[{AppleScript}] Copies to the clipboard the AppleScript code needed to render the
        selected shape; you can then paste the Applescript code into the Script Editor app and
        save that as an AppleScript.
      \item[{JavaScript}] Copies to the clipboard the JavaScript coded needed to render the
        selected shape; you can then paste the JavaScript into a text editor and save as an Omni
        Automation Plug-In.
    \end{description}
  \item[{Paste (\textbf{Command-V})}] Place the contents of the clipboard on the canvas, at the
    last place you clicked with the Selection tool. (Or insert text from the clipboard if you
    are editing an object's text.) If there is an image or some text in the clipboard, a new
    shape object containing the image or text is created.
  \item[{Paste in Place (\textbf{Shift-Command-V})}] If you copy an item on your canvas with
    \textbf{Command-C} and then use Paste in Place (\textbf{Shift-Command-V}), you will paste in
    the contents of the clipboard directly on top of the item you copied. This is particularly
    helpful when you want to copy an item on one layer and then paste in another copy on another
    layer that's higher up in the stack.
  \item[{Paste and Match Style (\textbf{Option-Shift-Command-V})}] When you are editing text,
    you can use this command to paste in some text without bringing along the style it had in
    the application you copied it from. Instead, it matches the style of the text you already
    have, as if you had typed it there yourself.
  \item[{Duplicate and Maintain Layers (\textbf{Command-D})}] Make a copy of the current
    selection and place it on the same layer. If you select objects on different layers and
    choose this option, duplicate copies of those objects will appear on their respective
    layers.

    \begin{ogpro}
      With an object selected, you can also hold down the \textbf{Option} and \textbf{Command} keys,
      and then press one of the four arrow keys (Up, Down, Left, or Right) to place a duplicate of
      the selected object in the direction of the arrow key you press. For example, pressing
      \textbf{Option-Command-Up Arrow} places the duplicate above the currently selected object.

      Objects are automatically placed 36 pixels apart. You can change the distance between
      objects using the Alignment inspector.
    \end{ogpro}
  \item[{Delete}] Obliterate the current selection, forever banishing it from this plane of
    existence (unless you Undo).
  \item[{Select All (\textbf{Command-A})}] Select everything selectable on the canvas. Remember
    that objects on locked layers can't be selected. If you are editing the text of an object,
    this instead selects all of the object's text.
  \item[{Select}] This menu offers you different options for selecting objects, either based on
    their style, how they're connected, or based on their relationships:

    \begin{description}[style=nextline]
      \item[{Similar Objects}] Select all of the objects on the canvas that have the same style
        attributes as the currently selected objects.
      \item[{Connected Objects}] Select all shapes that are directly or indirectly connected by
        lines to the currently selected objects, and all of the lines between.
      \item[{Ancestor}] Select all of the ancestors of the currently selected objects, based on
        the hierarchy established by connection lines.
      \item[{Descendants}] Select all of the descendants of the currently selected objects, based
        on the hierarchy established by connection lines.
    \end{description}
  \item[{Outlining}] These commands are available when you are working in the \textbf{Outline}
    tab of the Sidebar (\textbf{Option-Command--4}), or when focused on the canvas:

    \begin{itemize}
    \item Add Child (\textbf{Command-\}})---Adds an item nested inside the currently selected
      item.
    \item Add Sibling (\textbf{Option-Command-'})---Adds a peer of the currently selected item,
      represented on canvas by a shared relationship to the parent.
    \item Add Aunt (\textbf{Command-\{})---Adds an item that is a sibling of the selected item's
        parent.
      \item Add Parent (\textbf{Option-Command-;})---Adds an item at a level above the currently
        selected item, and moves the item (and any of its children) inside it.
      \item Add Spouse (\textbf{Option-Command-/})---Adds a peer of the currently selected item,
        represented on canvas by a shared relationship to the selected item.
      \item Indent---Moves the selected item one level lower in the outline.
      \item Outdent---Moves the selected item one level higher in the outline.
    \end{itemize}
  \item[{Navigation (\textbf{Standard})}] These commands are most useful when you invoke their
    keyboard shortcuts; you can do a whole lot without even touching the mouse:

    \begin{itemize}
      \item Select Shape Left (\textbf{Command-Left Arrow})
      \item Select Shape Right (\textbf{Command-Right Arrow})
      \item Select Shape Above (\textbf{Command-Up Arrow})
      \item Select Shape Below (\textbf{Command-Down Arrow})
    \end{itemize}
\end{description}

\begin{ogpro}
  \begin{description}[style=nextline]
    \item[{Mouseless Editing (\textbf{PRO})}] With OmniGraffle Pro, the previously mentioned
      Navigation commands are joined by options for creating a duplicate of the selected object
      in the direction of the arrow key you pressed. For example, with an object selected,
      pressing \textbf{Option-Command-Right Arrow} places a duplicate of the object on the
      canvas to the right. The spacing between objects is controlled by the \textbf{Alignment}
      inspector's horizontal and vertical offset values.

      \begin{itemize}
        \item Select Shape Left (\textbf{Command-Left Arrow})
        \item Select Shape Right (\textbf{Command-Right Arrow})
        \item Select Shape Above (\textbf{Command-Up Arrow})
        \item Select Shape Below (\textbf{Command-Down Arrow})
        \item (\textbf{PRO}) Create Shape Left (\textbf{Option-Command-Left Arrow})
        \item (\textbf{PRO}) Create Shape Right (\textbf{Option-Command-Right Arrow})
        \item (\textbf{PRO}) Create Shape Above (\textbf{Option-Command-Up Arrow})
        \item (\textbf{PRO}) Create Shape Below (\textbf{Option-Command-Down Arrow})
      \end{itemize}

      See Positioning Objects Neatly with the Alignment Inspector for more information about using
      the Alignment inspector.
  \end{description}
\end{ogpro}

\begin{description}[style=nextline]
  \item[{Canvases}] Commands related to working with the Canvas:

    \begin{description}[style=nextline]
      \item[{New Canvas (\textbf{Option-Command-N})}] Add a new canvas to the document.
      \item[{Duplicate Canvas}] Add a new canvas, identical to the current canvas, to the
        document.
      \item[{Delete Canvas}] Get rid of the current canvas completely.
    \end{description}
  \item[{Layers}] Commands related to working with Layers:

    \begin{description}[style=nextline]
      \item[{New Layer}] Add a new layer to the current canvas.
      \item[{New Layer From Selection}] Moves the selected object(s) to a new layer on the same
        canvas.
    \end{description}

    \begin{ogpro}
      \begin{description}[style=nextline]
        \item[{New Shared Layer (\textbf{PRO})}] Add a new shared layer to the current canvas.
        \item[{New Artboard Layer (\textbf{PRO})}] Add a new Artboard Layer to the current canvas.
      \end{description}
    \end{ogpro}

    \begin{description}[style=nextline]
      \item[{Duplicate Layer}] Add a new layer, identical to the selected layer, to the canvas.
      \item[{Merge Layer Down}] Move everything on the selected layer to the next layer down, then
        delete the selected layer.
      \item[{Delete Layer}] Get rid of the current layer completely. You can't delete the last
        layer in a canvas.
    \end{description}

    \begin{ogpro}
      \begin{description}[style=nextline]
        \item[{Delete from All Canvases (\textbf{PRO})}] This command is only available when you
          select a shared layer. It deletes all instances of the shared layer on any canvas
          throughout the document.
      \end{description}
    \end{ogpro}

    \begin{description}[style=nextline]
      \item[{Layer Settings}] Commands for controlling individual layers.

        \begin{description}[style=nextline]
          \item[{Visible}] shows or hides objects on the layer
          \item[{Printable}] controls whether objects on the layer should be included when printing
          \item[{Unlocked}] controls whether objects on the layer can be selected or edited
        \end{description}

        \begin{ogpro}
          \begin{description}[style=nextline]
            \item[{Shared Layer (\textbf{PRO})}] converts the current layer to a Shared Layer,
              making it accessible with other canvases in the project.
            \item[{Artboard Layer (\textbf{PRO})}] converts the current layer to an Artboard Layer.
          \end{description}
        \end{ogpro}
      \item[{Select All in Layer}] Select all of the objects on the selected layer.
      \item[{Move Selection to Layer}] Move all selected objects to the selected layer.
    \end{description}
\end{description}

\begin{ogpro}
  \begin{description}[style=nextline]
    \item[{Tables (\textbf{PRO})}] Commands for working with Tables:

      \begin{description}[style=nextline]
        \item[{Insert Row (\textbf{Command-Return})}] When you have one cell of a table selected,
          choose this command to add a new row of cells to the table. The new row appears
          immediately before the row containing the selected cell.
        \item[{Insert Column (\textbf{Option-Command-Return})}] When you have one cell of a table
          selected, choose this command to add a new column of cells to the table. The new row
          appears immediately before the row containing the selected cell.
        \item[{Select Row}] When you have one cell of a table selected, choose this command to
          select the entire row that contains the cell.
        \item[{Select Column}] When you have one cell of a table selected, choose this command to
          select the entire column that contains the cell.
      \end{description}
  \end{description}
\end{ogpro}

\begin{ogsidebar}
  Creating Tables (\textbf{PRO})

  A table is a special kind of group that organizes rows and columns of objects. To create a
  table:

  \begin{enumerate}
    \item Create a single shape object to be the first cell of your table.

      \ogfigure{0.704}{res/og7_tables01.png}
    \item With the object selected, choose \textbf{Arrange} $\blacktriangleright$ \textbf{Make
      Table} (\textbf{Shift-Command-T}). Notice the grid-like handles on the top, bottom, left,
      and right sides of the object; this lets you know the object is a table.

      \ogfigure{0.704}{res/og7_tables02.png}
    \item Drag the right-side table handle to the right to add columns.

      \ogfigure{0.800}{res/og7_tables03.png}
    \item Drag the bottom table handle down to add rows.

      \ogfigure{0.800}{res/og7_tables04.png}
  \end{enumerate}

  You can also select a number of already-existing objects and then choose \textbf{Arrange}
  $\blacktriangleright$ \textbf{Make Table}; OmniGraffle rearranges the objects into rows and
  columns to the best of its ability.

  \ogfigure{0.800}{res/og7_tables05.png}

  Table cells can be selected and edited just like group members. If you resize a cell, the rest
  of the cells in its row and column resize accordingly. When selecting cells, you can
  Shift-click to select everything between the cell you click and the last cell you selected.

  You can add or remove rows or columns at any time by dragging the table handles outward or
  inward. Resize the table as a whole by dragging the ordinary corner handles. Use the
  \textbf{Edit} $\blacktriangleright$ \textbf{Tables} commands when working within table cells
  to insert or select rows and cells.

  To dismantle a table, select it and choose \textbf{Arrange} $\blacktriangleright$
  \textbf{Ungroup} (\textbf{Shift-Command-U}).
\end{ogsidebar}

\begin{description}[style=nextline]
  \item[{Magnets}] Commands for controlling the Magnets attached to objects:

    \begin{description}[style=nextline]
      \item[{Copy Magnets}] Put the magnet arrangement of the selected object onto the clipboard,
        so that you can paste them to a different object.
      \item[{Paste Magnets}] If there is a magnet arrangement on the clipboard, apply it to the
        selected object. Any existing magnets are replaced by the new ones.
      \item[{Delete Magnets}] Get rid of all the magnets on the selected objects.
      \item[{Positions}] In addition to these options, magnet positions can also be set in the
        \textbf{Connections} properties inspector (\textbf{Command--3}) in the \textbf{Other
        Connections} tab:
    \end{description}

    \begin{itemize}
      \item No magnets --- removes all Magnets from the selected shape
      \item Cardinal Points:

        \begin{itemize}
          \item 4 magnets: N, S, E, W
          \item 2 magnets: N, S
          \item 2 magnets: E, W
          \item 4 magnets: NE, NW, SE, SW
          \item 8 magnets
        \end{itemize}
      \item Per Segment:

        \begin{itemize}
          \item on each vertex --- makes it so each point on a shape's path is also a Magnet
          \item 1 magnet per side
          \item 2 magnets per side
          \item 3 magnets per side
          \item 4 magnets per side
          \item 5 magnets per side
        \end{itemize}
    \end{itemize}
\end{description}

\begin{ogpro}
  \begin{description}[style=nextline]
    \item[{Objects (\textbf{PRO})}] Commands for working with Shapes:

      \begin{description}[style=nextline]
        \item[{Unite Shapes}] Combine two or more selected shapes, eliminating internal lines to
          make a single, solid shape.
        \item[{Intersect Shapes}] Create a shape from the overlapping areas of the selected
          shapes.
        \item[{Subtract Shapes}] Take the selected shapes and use the front shapes to punch a hole
          in the back-most shape. You can reorder shapes with the Bring/Send commands in the
          Arrange menu.
        \item[{Uncombine Shapes}] Separates combined shapes, restoring them to their original
          state.
        \item[{Convert Text to Shapes}] Converts the selected text or text object to an editable
          shape.
        \item[{Convert Line to Shape}] Converts the selected line or stroke to an editable shape.
        \item[{Remove Midpoints and Reset Handles}] Removes any midpoints added to the object, and
          resets the position of its control handles.
        \item[{Flip Horizontally}] Flips the selected shape, left-to-right.
        \item[{Flip Vertically}] Flips the selected shape, top-to-bottom.
      \end{description}
  \end{description}
\end{ogpro}

\begin{ogsidebar}
  Text to Shapes (\textbf{PRO})

  An age-old problem for designers is the ability to do more with the text they're working with.
  Whether it's adding more style to the text by adding a different fill or making slight
  alterations to how a character appears, can sometimes be really, really hard.

  With OmniGraffle 7, though, we've taken the pain out of this task, giving you the ability to
  convert any text object to shapes. Once converted, you can do way more with the text shapes
  than you could with a standard text element in previous versions of OmniGraffle.

  \subsection{Converting Text to Shapes}

  To convert text to shapes, you'll first need a text object on the Canvas. Use the
  \textbf{Text} tool in combination with the \textbf{Font} and \textbf{Text Position}
  inspectors (see Working With the Inspectors, earlier in this guide) to set the text
  properties prior to clicking on the Canvas. Then, with the text object selected, choose
  \textbf{Edit} $\blacktriangleright$ \textbf{Objects} $\blacktriangleright$ \textbf{Convert
  Text to Shapes}.

  \ogfigure{0.800}{res/og7_texttoshapes01.png}

  OmniGraffle automatically groups the shaped-letters together:

  \ogfigure{0.800}{res/og7_texttoshapes02.png}

  To work on the letters individually, you can select the individual shapes and then use the
  \textbf{Point Editor} tool to change the appearance of the characters. For example, by
  selecting one of the blue vector points you could:

  \begin{itemize}
    \item delete a vector point by selecting the point and pressing the \textbf{Delete}
      key\ldots{}

      \ogfigure{0.800}{res/og7_texttoshapes04.png}
    \item \ldots{}change the Bézier curve by using the control handles\ldots{}

      \ogfigure{0.800}{res/og7_texttoshapes05.png}
    \item \ldots{}or drag and move the vector points around.

      \ogfigure{0.800}{res/og7_texttoshapes06.png}
  \end{itemize}

  \subsection{Uses for Text to Shapes}

  With a standard text object, all you can do to that text---style-wise---is apply a single fill
  color, maybe change its opacity, and apply a drop shadow. Pretty boring stuff, actually. If
  you try to apply a fill color or pattern, or change the stroke of the text, the changes you
  make in the Fill and Stroke inspectors are applied to the shape of the text object, not the
  text.

  However, after converting text to shapes, you can use all of the available options in the Fill
  and Stroke inspectors to your delight! For example, you could apply a gradient fill with a
  beveled stroke:

  \ogfigure{0.800}{res/og7_texttoshapes07.png}

  Or you could select a \textbf{Fill Blend Type} and \textbf{Distortion Effect} in the Fill
  inspector, and apply that to text on a layer above an image or some other object:

  \ogfigure{0.800}{res/og7_texttoshapes08.png}

  And, since each character is now a shape, you can use the \textbf{Shape} inspector to change
  any character into another shape!

  \ogfigure{0.800}{res/og7_texttoshapes09.png}

  One more benefit of converting text to shapes is to allow for more precise kerning control.
  You can select each character individually, and nudge them left or right by pressing the
  \textbf{Left} or \textbf{Right Arrow} keys.

  So, go ahead, let your creative side run wild with OmniGraffle 7's new \textbf{Text to
  Shapes}.
\end{ogsidebar}

\begin{ogsidebar}
  Lines to Shapes (\textbf{PRO})

  To give designers more flexibility with their work, we have added the ability convert the
  lines you draw with the \textbf{Line} tool into paths. After converting a line to a path,
  you can use the \textbf{Point Editor} tool to add or remove vector points from the path, add
  Bézier curves, and click and drag vector points to ``open'' the path.

  Converting Lines to Shapes

  Converting a line to a shape is fairly simple:

  \begin{enumerate}
    \item First, draw a line on the canvas with the \textbf{Line Tool}. If you click-and-drag
      when creating a vector point, you instantly convert that point into a Bézier curve:

      \ogfigure{0.800}{res/og7_linetoshape01.png}
    \item Choose \textbf{Edit} $\blacktriangleright$ \textbf{Objects} $\blacktriangleright$
      \textbf{Convert Line to Shape}

      \ogfigure{0.800}{res/og7_linetoshape02.png}
  \end{enumerate}

  Once converted, the ``line'' shape appears on the canvas as a shape with an empty fill.

  Working with Line Shapes

  If you select the \textbf{Point Editor} tool and zoom in on the shape, you will see the blue
  diamond-shaped vector points on the path that make up the shape.

  \ogfigure{0.800}{res/og7_linetoshape03.png}

  You can use the Point Editor tool to:

  \begin{itemize}
    \item Edit Bézier curves by clicking and dragging the red control handles
    \item Add vector points to the path by double-clicking on the path
    \item Delete vector points by first selecting the point and then pressing \textbf{Delete}
  \end{itemize}

  Similarly, you can use the Inspectors to:

  \begin{itemize}
    \item Add a Fill to the shape
    \item Change or remove the shape's Stroke
    \item Add a Shadow
    \item Fill the shape with text or an image
  \end{itemize}

  \ogfigure{0.800}{res/og7_linetoshape04.png}

  Uses for Line Shapes

  Lines on their own can make a dramatic impact on your designs. However, the ability to convert
  them to shapes and apply various effects opens the door to a realm of possibilities.

  For example, you can fill the shape with a gradient and remove its stroke to create gradient
  rule lines, or empty lines with inner shadows:

  \ogfigure{0.800}{res/og7_linetoshape05.png}

  Use the \textbf{Point Editor} tool to remove vector points from the paths, add new points, and
  conver them to Béziers:

  \ogfigure{0.800}{res/og7_linetoshape06.png}

  Use the \textbf{Point Editor} tool to open a line shape, and then place an image into the
  shape to create a custom image \textit{mask}:

  \ogfigure{0.800}{res/og7_linetoshape07.png}
\end{ogsidebar}

\begin{description}[style=nextline]
  \item[{Mask Image/End Masking}] Resize or rotate an image object on the canvas. Hold
    \textbf{Shift} to scale the image, and hold \textbf{Command} while hovering over one of the
    image handles to rotate an image. Holding \textbf{Shift-Command} lets you rotate the image
    in 15 degree increments. Press \textbf{Return}, or choose \textbf{Edit}
    $\blacktriangleright$ \textbf{End Masking} when finished to accept any masking changes
    you've made to the image.
\end{description}

\begin{ogtip}
  While not exactly the same as cropping an image, using an image mask is a great way to get a
  crop effect by making only a selected portion of the image visible.
\end{ogtip}

\begin{description}[style=nextline]
  \item[{Insert Variable}] While you are editing a shape object's text, choose one of these
    commands to insert a special tag. When you finish editing the text, the tag changes to
    represent some aspect of the object itself, the canvas, the page, or the document. If the
    data that the tag refers to changes, the tag automatically updates itself. Some variables,
    such as Page Number, are useful for putting on a shared layer and then sharing that layer
    with every canvas.

    \begin{description}[style=nextline]
      \item[{Canvas Name \texttt{<\%Canvas\%>}}] the name of the canvas the object inhabits.
      \item[{Layer Name \texttt{<\%Layer\%>}}] the name of the layer.
      \item[{Document Name \texttt{<\%Document\%>}}] the name of the document the object inhabits.
      \item[{Current Date \texttt{<\%Date\%>}}] just what it says; this option inserts the current
        date. This particularly handy when you want to ``stamp'' a document you're printing with
        today's date.
      \item[{Document Creation Date \texttt{<\%CreationDate\%>}}] the date that the document was
        originally created, according to macOS's filesystem.
      \item[{Document Modification Date \texttt{<\%ModificationDate\%>}}] the date that the
        document was last saved, according to macOS's filesystem.
    \end{description}

    \begin{ogsidebar}
      Understanding Date Variables

      If you are using one of the date variables, you can customize the format of the date using
      the variables listed here:

      \begingroup
      \small
      \setlength{\LTleft}{0.075\textwidth}
      \setlength{\LTright}{0.075\textwidth}
      \renewcommand{\arraystretch}{1.25}
      \begin{xltabular}{0.85\textwidth}{|l|>{\raggedright\arraybackslash}X|}
        \hline
        \rowcolor{ogheaderblue}
        \textcolor{white}{\textbf{Variable}} & \textcolor{white}{\textbf{Description}} \\
        \hline
        \endhead
        \hline
        \endfoot
        \ogcellcode{y} & Two-digit year, based on the decade; for example, 16 for ``2016'' \\
        \hline
        \ogcellcode{Y} & Four-digit year; for example, 2016 \\
        \hline
        \ogcellcode{m} & Month, represented as two digits ranging from 01-12; for example, 10
        for October \\
        \hline
        \ogcellcode{d} & Day, represented as two digits ranging from 01-31 (as appropriate); for
        example, 05 is the fifth day of the month \\
        \hline
        \ogcellcode{H} & Hours, represented as two-digits ranging from 00-23; for example, 11 is
        11 a.m., whereas 23 is 11 p.m. \\
        \hline
        \ogcellcode{M} & Minutes, represented as two digits ranging from 00-59 \\
        \hline
        \ogcellcode{S} & Seconds, represented as two digits ranging from 00-59 \\
        \hline
        \ogcellcode{z} & Zulu time difference; for example, an output of -0700 represents that
        the current time is seven hours behind Zulu time \\
        \hline
      \end{xltabular}
      \endgroup

      For example, to get a date such as \texttt{2016-10-05 08:37:30 -0700Z}, you could use the
      following variable string:

      \texttt{<\%date \%Y-\%m-\%d \%H:\%M:\%S \%zZ \%>}

      The logic we have used in creating the date variables is based on
      \href{http://developer.apple.com/library/ios/documentation/cocoa/Conceptual/DataFormatting/Articles/dfDateFormatting10_4.html}{Cocoa's
      date formatting tokens}.
    \end{ogsidebar}

    \begin{description}[style=nextline]
      \item[{Creator \texttt{<\%Creator\%>}}] the name of the user who created the file, according
        to macOS's filesystem.
      \item[{Modifier \texttt{<\%Modifier\%>}}] the name of the user who last saved the file,
        according to macOS's filesystem.
      \item[{Page Number \texttt{<\%\#\%>}}] the number of the page that the object inhabits,
        based on settings in the Canvas Size inspector and the Page Setup panel.
      \item[{Total Pages \texttt{<\%TotalPages\%>}}] the number of pages in the whole document.
        Use this with the Page Number variable to get results like ``Page 2 of 5''; for example:
      \item[{}] \texttt{<\%\# of \%TotalPages \%>}
      \item[{Canvas Number \texttt{<\%CanvasNumber\%>}}] if you have multiple canvases in your
        project, you can reference them by number.
      \item[{Total Canvases \texttt{<\%TotalCanvases\%>}}] prints the number of canvases in the
        document. Use this with the Canvas Number variable to get results like ``Canvas 1 of 3'';
        for example:
      \item[{}] \texttt{<\%CanvasNumber of \%TotalCanvases \%>}
      \item[{X Position \texttt{<\%X\%>}}] the object's horizontal distance from the origin, just
        like in the Geometry inspector.
      \item[{Y Position \texttt{<\%Y\%>}}] the object's vertical distance from the origin, just
        like in the Geometry inspector.
      \item[{Width \texttt{<\%Width\%>}}] the object's horizontal size, just like in the Geometry
        inspector.
      \item[{Height \texttt{<\%Height\%>}}] the object's vertical size, just like in the Geometry
        inspector.
      \item[{Line Length \texttt{<\%Length\%>}}] Apply the \textbf{Line Length} variable to the
        label on a line. The Line Length variable relies on the values set in the \textbf{Units}
        inspector when displaying the line length. For example, let's say that you are creating a
        floorplan of your house or apartment. In the \textbf{Units} inspector, you can set the
        \textbf{Ruler Units} to \textbf{feet (ft)}, and the \textbf{Scale} to something like
        \textbf{1 in = 1 ft} to get accurate (to-scale) floorplans, using lines and the Line
        Length variable to display the dimensions.

        \ogfigure{0.769}{res/og7_menus01.png}
      \item[{GraphicName \texttt{<\%GraphicName\%>}}] Places on the canvas a text object with the
        name of the currently selected shape.
    \end{description}
  \item[{Find}] Commands for searching within OmniGraffle:

    \begin{description}[style=nextline]
      \item[{Find (\textbf{Command-F})}] This summons OmniGraffle's Find panel. You can search for
        a specific string of text, or you can use a regular expression to match a whole class of
        strings. When searching for a specific string, select \textbf{Ignore Case} to match
        letters regardless of whether they are uppercase or lowercase; select Whole Word to avoid
        matching portions of a word. If you use groups in a regular expression, you might want to
        select (or replace) the whole expression, or just one of the groups; use the pop-up menu
        to indicate which.
    \end{description}

    \begin{ogsidebar}
      Searching with Regular Expressions

      The Find dialog supports the use of
      \href{https://en.wikipedia.org/wiki/Regular_expression}{regular expressions} to match more
      complex text strings and combinations. See the
      \href{https://developer.apple.com/documentation/foundation/nsregularexpression}{NSRegularExpression
      developer documentation} for technical details on the regex syntax, including a list of
      supported characters and operators.
    \end{ogsidebar}

    \begin{description}[style=nextline]
      \item[{Find Next (\textbf{Command-G})}] Choose this to find the next occurrence of the text
        you last entered in the Find panel, as if you had just clicked its Next button.
      \item[{Find Previous (\textbf{Shift-Command-D})}] Choose this to find the previous
        occurrence of the string you last entered in the Find panel, as if you had just clicked
        its Previous button.
    \end{description}

    \begin{ognote}
      For the purposes of Previous and Next, OmniGraffle considers objects in left-to-right,
      top-to-bottom order.
    \end{ognote}

    \begin{description}[style=nextline]
      \item[{Enter Selection (\textbf{Command-E})}] Select some text in an object, then choose
        this command to put the selected text into the Find panel, replacing whatever was there
        before. Then you can use the Find Next and Find Previous commands to search for other
        occurrences of the text.
      \item[{Scroll to Selection (\textbf{Command-J})}] Move the view so that the selected objects
        are visible.
      \item[{Zoom to Selection}] Fits the currently selected object(s) so that they fit within the
        Canvas view.
    \end{description}
  \item[{Spelling and Grammar}] Commands for checking the spelling and grammar within your
    project file:

    \begin{description}[style=nextline]
      \item[{Show Spelling and Grammar (\textbf{Command-:})}] Opens the Spelling and Grammar
        dialog. Since the colon (\textbf{:}) is on the upper part of semicolon (\textbf{;}) key,
        you will also need to use the Shift key to option this dialog (\textbf{Shift-Command-;}).
      \item[{Check Document Now (\textbf{Command-;})}] Checks the spelling in your current
        document and stops at the first misspelled word it finds.
      \item[{Check Spelling While Typing}] This option is enabled by default, but can be turned
        off by selecting this option in the menu.
      \item[{Check Grammar With Spelling}] This option also checks your grammar as it checks
        spelling. This option comes in particularly handy when you're entering a lot of text.
    \end{description}

    \begin{ognote}
      For the purposes of Previous and Next, OmniGraffle considers objects in left-to-right,
      top-to-bottom order.
    \end{ognote}
  \item[{Transformations}] Commands for transforming a text selection:

    \begin{description}[style=nextline]
      \item[{Make Lower Case}] Replaces all upper case characters in the selected text with their
        lower case equivalents.
      \item[{Make Upper Case}] Replaces all lower case characters in the selected text with their
        upper case equivalents.
      \item[{Capitalize}] Converts the first letter of each word in the selection to its
        capitalized equivalent.
    \end{description}
  \item[{Edit LinkBack Item}] If you have pasted LinkBack content from another app, select the
    object and use this command to open the content in the original app.
  \item[{Start Dictation (\textbf{fn-fn})}] To use Dictation, place a text object on the canvas
    or double-click a shape to make its its label field active, and then press the \textbf{fn}
    key twice to start dictation. The words you speak are inserted as text. To stop Dictation,
    press \textbf{fn} twice.
  \item[{Emoji \& Symbols (\textbf{Control-Command-Space})}] Summon the standard macOS Character
    Palette.
\end{description}

\section{The View Menu}

The options in the View menu provide ways for you to zoom in or out, and control the way you
look at the Canvas.

\begin{description}[style=nextline]
  \item[{Show/Hide Tab Bar}] If you are using macOS Sierra (v10.12, or higher), you can enable a
    tab bar within OmniGraffle so you can work on multiple documents. With multiple tabs open,
    you can use the tab-related options in the Window menu to switch between tabs, move a tab to
    another window, or to collect all OmniGraffle windows as tabs in the current document
    window.
\end{description}

\begin{ogpro}
  \begin{description}[style=nextline]
    \item[{Start Presentation (\textbf{PRO}) (\textbf{Option-Command-P})}] Enter presentation
      mode with the frontmost document.
  \end{description}
\end{ogpro}

\begin{description}[style=nextline]
  \item[{Zoom}] These menu options let you zoom in and out on the canvas, and are helpful for
    when you need fine-grained control over where objects are placed:

    \begin{description}[style=nextline]
      \item[{Zoom In (\textbf{Command->})}] Move up to the next closest zoom level, making the
        canvas appear larger.
      \item[{Zoom Out (\textbf{Command-<})}] Move down to the next farthest zoom level, making the
        canvas appear smaller.
      \item[{Zoom to Actual Size (\textbf{Option-Command--0})}] Move to 100 percent zoom, so that
        the canvas appears as large as it actually is according to the canvas units.
      \item[{Zoom to Selection (\textbf{Option-Command-J})}] Zoom to a level and position that
        fills the view with the currently selected objects.
      \item[{Previous Zoom}] Return to the previous zoom level you had before the current one.
      \item[{Fit in Window}] Zoom to a level and position that keep every object on the canvas
        visible as you work.
    \end{description}
  \item[{Extras}] There are several things on the canvas (other than the objects themselves)
    that you might want to see; choose something from this list to show or hide it.

    \begin{description}[style=nextline]
      \item[{Actions}] Any object with an action assigned to it in the Action inspector gets an
        action tag near its upper-right corner.
      \item[{Grid Lines (\textbf{Command-\textbackslash\{\}})}] This is the grid as defined in the
        Grid inspector. Note that you can snap to the grid whether it's visible or not.
      \item[{Guide Lines}] These are the manual guides that you set up for aligning objects on the
        canvas.
      \item[{Magnets}] Choose this to see all objects' magnets, all the time.
    \end{description}

    \begin{ogpro}
      \begin{description}[style=nextline]
        \item[{Notes (\textbf{PRO})}] Any object with a note in the Object Data inspector gets a
          note tag near its upper-right corner.
      \end{description}
    \end{ogpro}

    \begin{description}[style=nextline]
      \item[{Origin}] This is the crosshair that represents the canvas origin.
      \item[{Page Breaks}] These are the lines that represent borders between pages.
      \item[{Rulers (\textbf{Command-R})}] The rulers help you size objects on the canvas, and
        include an inspector bar for common object editing.
      \item[{Movement Handles}] Choose whether to have object movement handles displayed. When
        displayed, a handle is placed at the center of the object, which you can click and drag to
        move the object; clicking and dragging the other handles will resize the object in the
        direction you drag the mouse.
    \end{description}
  \item[{Show/Hide Sidebar (\textbf{Option-Command--1})}] Toggle the visibility of the left
    sidebar which contains Canvases and their layers, and Contents.
  \item[{Show Contents}] Use these options to toggle the Sidebar's view options:

    \begin{description}[style=nextline]
      \item[{Layers (\textbf{Option-Command--2})}] This is the default view for the Sidebar, which
        shows you the canvases, layers, and all associated objects on those layers.
      \item[{Guides (\textbf{Option-Command--3})}] Displays the positions and colors of the
        Guides.
      \item[{Outline (\textbf{Option-Command--4})}] The Outline Editor makes it easy for you to
        quickly enter---and edit---objects in a project. This provides you with a hierarchical
        view of the objects in your project so you can quickly assess their relationships.
      \item[{Selection (\textbf{Option-Command--5})}] The Selection Matrix groups similarly-styled
        objects---blocks, labels, lines, images, together so you can quickly select objects based
        on their style.
    \end{description}
  \item[{Display Canvas}] These options make it easy for you to switch between the different
    canvases in your project:

    \begin{description}[style=nextline]
    \item[{Next Canvas (\textbf{Command-]})}] Switch to the canvas after the one you are
      currently editing.
    \item[{Previous Canvas (\textbf{Command-[})}] Switch to the canvas before the one you are
        currently editing.
    \end{description}
  \item[{Customize Toolbar}] Conjure up a sheet of controls that you can drag to your toolbar,
    as in other Mac apps.
  \item[{Show/Hide Toolbar}] Set the visibility of the document window's toolbar.
  \item[{Enter/Exit Full Screen (\textbf{Control-Command-F})}] Enter full screen mode.
\end{description}

\section{The Format Menu}

The options in the Format menu provide control over fonts, colors, and object and diagram
styles.

\begin{description}[style=nextline]
  \item[{Show Fonts (\textbf{Command-T})}] Opens the standard Font panel.
  \item[{Show Colors (\textbf{Shift-Command-C})}] Opens the standard Color panel.
\end{description}

\begin{ogsidebar}
  Using the Color Picker

  OmniGraffle uses the standard macOS color picker for selecting the colors that you use
  throughout the app.

  \ogfigure{0.800}{res/og7_colorpicker.png}

  The color picker includes the following features:

  \begin{itemize}
    \item \textbf{Color Selection Inspectors} --- The color picker offers several modes for
      inspecting and choosing colors for your document: a Color Wheel, Color Sliders, Color
      Palettes, Image Palettes, Pencils, and Pattern Palettes. Use the tabs at the top of the
      panel to switch between them.
    \item \textbf{Opacity Slider} --- Beneath the inspector there is a slider for setting the
      opacity of the selected color, from 0\% (transparent) to 100\% (fully opaque). The color
      well in the lower left reflects the difference between the opacity setting and the
      original color.
    \item \textbf{Eyedropper Tool} --- Next to the color well in the lower left is the
      eyedropper tool. Click it and then click anywhere on the screen to select the color that
      appears at that location.
    \item \textbf{Color Chits} --- The lower right section of the color picker contains rows of
      color chits, customizable swatches that you can fill with colors of your choice by
      dragging from the color well to the left. To reveal more chits, expand the panel by
      dragging its lower right corner, or drag the dividing line between the chits and the
      opacity slider.

      To remove a custom chit from the color picker, drag it to the macOS Trash in the Dock.
  \end{itemize}
\end{ogsidebar}

\begin{description}[style=nextline]
  \item[{Font}] Use this submenu to set a font and style to the selected text.

    \begin{description}[style=nextline]
      \item[{Bold (\textbf{Command-B})}] Make the selected text \textbf{bold}.
      \item[{Italic (\textbf{Command-I})}] Make the selected text \textit{italic}.
      \item[{Underline (\textbf{Command-U})}] \underline{Underline} the selected text.
      \item[{Outline}] Places an outline around the selected text.
      \item[{Strikethrough}] Adds a strikethrough line to the selected text.
      \item[{Bigger (\textbf{Command-+})}] Make the selected text bigger.
      \item[{Smaller (\textbf{Command---})}] Make the selected text smaller.
      \item[{Kern}] Change the spacing between characters; options include:

        \begin{itemize}
          \item Use Default
          \item Use None
          \item Tighten
          \item Loosen
        \end{itemize}
      \item[{Ligature}] If the font of the selected text offers ligatures, you can use these
        options to join two or more characters into a single glyph:

        \begin{itemize}
          \item \textbf{Use Default} --- characters appear with minimal ligatures, if available.
        \end{itemize}

        \ogfigure{0.481}{res/og7_ligature_default.png}

        \begin{itemize}
          \item \textbf{Use None} --- characters appear without ligatures.
        \end{itemize}

        \ogfigure{0.481}{res/og7_ligature_none.png}

        \begin{itemize}
          \item \textbf{Use All} --- characters are joined into a single glyph, if available.
        \end{itemize}

        \ogfigure{0.481}{res/og7_ligature_all.png}
    \end{description}

    \begin{ognote}
      For more information about ligatures, see the Wikipedia page on
      \href{https://en.wikipedia.org/wiki/Typographic_ligature}{Typographic ligatures}.
    \end{ognote}

    \begin{description}[style=nextline]
      \item[{Baseline}] Alter the text's baseline; options include:

        \begin{itemize}
          \item Use Default
          \item Superscript
          \item Subscript
          \item Raise
          \item Lower
        \end{itemize}
    \end{description}
  \item[{Copy Style (\textbf{Option-Command-C})}] Copies the style attributes of the selected
    object.
  \item[{Paste Style (\textbf{Option-Command-V})}] Pastes the style attributes to the selected
    object.
  \item[{Text}] These options give you control over the selected text's properties, including:

    \begin{itemize}
      \item Align Left (\textbf{Control-Command-\{})
        \item Center (\textbf{Control-Command-|})
        \item Justify
      \item Align Right (\textbf{Control-Command-\}})
      \item Align Top (\textbf{Option-Command-\{})
        \item Center Vertically (\textbf{Option-Command-|})
      \item Align Bottom (\textbf{Option-Command-\}})
      \item Center Horizontally \& Vertically
      \item Allow Hyphenation
    \end{itemize}
  \item[{Copy Object Style (\textbf{Control-Command-C})}] Put the styles of the selected object
    on the clipboard, where they can be pasted onto other objects.
  \item[{Paste Object Style (\textbf{Control-Command-V})}] If there are object styles on the
    clipboard, apply them to the selected objects.
\end{description}

\section{The Arrange Menu}

Use these commands for working with objects on the Canvas or within Layers.

\begin{description}[style=nextline]
  \item[{Bring to Front (\textbf{Shift-Command-F})\\ Bring Forward (\textbf{Option-Command-F})}]
    Objects on the same layer have an order to them: newly created objects appear in front of
    objects that already existed. Choose Bring to Front to move the selected objects in front of
    all other objects on the layer. Hold Option to change this command to Bring Forward, which
    moves the selected objects forward one step in the ordering.
  \item[{Send to Back (\textbf{Shift-Command-B})\\ Send Backward (\textbf{Option-Command-B})}]
    Similar to \textbf{Bring to Front}, except it moves the selected objects behind all other
    objects on the layer. Hold Option to change this command to Send Backward, which moves the
    selected objects backward one step in the ordering.
  \item[{Lock (\textbf{Command-L})}] Lock the selected objects. A locked object can't be moved
    or edited. When you try to select a locked object, it shows X marks instead of selection
    handles. Basically, the only thing you can do with a locked object is unlock it; the
    exception is that you can change an object's label and hierarchical position in the outline
    view. Locking is useful for making sure you don't accidentally change something you have set
    up just right.
  \item[{Unlock (\textbf{Option-Command-L})}] Unlock the selected objects so that they can be
    moved or edited.
  \item[{Group (\textbf{Shift-Command-G})}] Use this command with at least two objects selected
    to turn them into a group, making them behave like a single object.
  \item[{Ungroup (\textbf{Shift-Command-U})}] Use this command with a group selected to separate
    it back into its component objects.
\end{description}

\begin{ogpro}
  \begin{description}[style=nextline]
    \item[{Make Table (\textbf{PRO}) (\textbf{Shift-Command-T})}] Use this command with one
      object selected to turn the object into a 1 $\times$ 1 table, to which you can then add
      rows and columns.
    \item[{Make Artboard (\textbf{PRO}) (\textbf{Shift-Command-A})}] Use this command to wrap
      the selected objects in an artboard. For more details, see Using Artboards (PRO), earlier
      in this guide.
    \item[{Group as Subgraph (\textbf{PRO})}] Turn the selected objects into a self-contained
      subgraph that you can subsequently collapse to a single object.
    \item[{Collapse Subgraph (\textbf{PRO})}] Change the selected subgraph into a single object,
      retaining all of the connection lines coming in or out of it.
    \item[{Expand Subgraph (\textbf{PRO})}] Return a collapsed subgraph to its fully expanded
      majesty.
  \end{description}
\end{ogpro}

\begin{ogsidebar}
  Using Subgraphs (\textbf{PRO})

  A subgraph is a special kind of group that can be expanded to show the hierarchy inside it, or
  collapsed to make it a single compact object. To create a subgraph:

  \begin{enumerate}
    \item Select some objects on the canvas:

      \ogfigure{0.800}{res/og7_subgraphs01.png}
    \item With the objects selected, choose \textbf{Arrange} $\blacktriangleright$ \textbf{Group
      as Subgraph}. Your objects are enclosed by a subgraph rectangle:

      \ogfigure{0.800}{res/og7_subgraphs02.png}
    \item Style and label the objects inside the subgraph however you like:

      \ogfigure{0.800}{res/og7_subgraphs03.png}
  \end{enumerate}

  To collapse a subgraph, select it and then choose \textbf{Arrange} $\blacktriangleright$
  \textbf{Collapse Subgraph} (or from the contextual menu that appears when you Control-click
  the subgraph). The objects contained in the subgraph are hidden away, and the subgraph
  becomes a single object. Any incoming or outgoing connection lines belonging to objects in
  the subgraph become incoming or outgoing lines of the collapsed object itself.

  \ogfigure{0.800}{res/og7_subgraphs04.png}

  Subgraph members can be selected and edited just like group members: click once to select the
  subgraph, then click again to select a member inside. See Grouping and Ungrouping Objects
  for details on how to view, add to, or remove a Subgroup's objects in the Sidebar.

  You can resize the containing rectangle, as long as you don't make it too small to contain its
  member objects.

  To disassemble a subgraph, select it and choose \textbf{Arrange} $\blacktriangleright$
  \textbf{Ungroup} (\textbf{Shift-Command-U}).
\end{ogsidebar}

\begin{description}[style=nextline]
  \item[{Guides}] OmniGraffle offers Smart Guides and manual guides. Smart guides are little
    blue lines that appear as you move, resize, or rotate objects, to help you keep things
    aligned and spaced evenly.
  \item[{}] Manual guides are canvas-spanning magenta lines that you place yourself; you can
    make as many as you like, use them to align objects, and show or hide them from the View
    menu.

    \begin{description}[style=nextline]
      \item[{Smart Alignment Guides}] These are the Smart Guides that appear to help you line up
        objects' edges or centers. They also help you match objects' rotation when you
        Command-drag an object's selection handles.
      \item[{Smart Distance Guides}] These are the Smart Guides that help you space out objects
        evenly.
      \item[{Lock/Unlock Guides}] As long as the manual guides are locked, you won't be able to
        move them.
      \item[{New Horizontal/Vertical Guide}] This puts a guide across the canvas; if the guides
        are not locked you can use the Selection tool to drag it into position wherever you like.
        As you drag an object near a guide, its edges or center snap into place along the guide.
        You can also drag new guides out from the vertical and horizontal rulers.
      \item[{Clear Manual Guides}] Remove all of the manual guides from the canvas. If you just
        want to temporarily hide them, you can do so from the View menu.
    \end{description}
  \item[{Align}] These commands replicate the functionality of buttons in the Alignment
    inspector (\textbf{Command--1}), lining up objects' edges and centers, or spacing objects
    out evenly. Options include:

    \begin{itemize}
      \item Align Left Edges
      \item Align Right Edges
      \item Align Top Edges
      \item Align Bottom Edges
      \item Align Vertical Centers
      \item Align Horizontal Centers
      \item Make Centered Row
      \item Make Centered Column
      \item Align Edges to Grid
      \item Align Centers to Grid
      \item Align to Canvas Left Edge
      \item Align to Canvas Right Edge
      \item Align to Canvas Top Edge
      \item Align to Canvas Bottom Edge
      \item Align to Canvas Vertical Center
      \item Align to Canvas Horizontal Center
      \item Make Centered Row to Canvas
      \item Make Centered Column to Canvas
      \item Offset Vertically by \textit{X}px\textsuperscript{*}
      \item Offset Horizontally by \textit{X}px\textsuperscript{*}
    \end{itemize}

    \begin{ognote}
      The \textit{X}px\textsuperscript{*} distances associated with the \textbf{Offset Vertically
      by} and \textbf{Offset Horizontally by} options are determined by the distances set in the
      Alignment inspector. By default, these are set to 36 pixels.
    \end{ognote}
  \item[{Grid}] Snap to Grid and Align Objects to Grid work just like they do on the Grid
    inspector; Size Objects to Grid resizes objects to make their edges follow the nearest grid
    lines. Options include:

    \begin{itemize}
      \item Snap to Grid
      \item Align Objects to Grid (\textbf{Option-Command-[})
      \end{itemize}
    \item[{Diagram Layout}] These commands lay out objects, either all objects on the canvas or
      the selected objects, according to the settings in the Diagram Layout and Style inspector.
      Options include:

      \begin{itemize}
        \item Lay Out Entire Canvas
        \item Lay Out Selection
        \item Object Rank

          \begin{itemize}
            \item Default
            \item Minimum
            \item Maximum
            \item Same
          \end{itemize}
      \end{itemize}
    \item[{Size}] If you have objects of different sizes on the canvas, but would like them to be
      the same size, first select the object that has the dimensions you'd like the others to
      emulate, choose \textbf{Arrange} $\blacktriangleright$ \textbf{Size}, and then choose one of
      the resizing options:

      \begin{itemize}
        \item Make Same Width
        \item Make Same Height
        \item Make Same Size
        \item Make Natural Size
        \item Size to Fit Image
      \end{itemize}
    \item[{}] Make Natural Size makes an object's height and width equal. For objects with an
      image placed in the Image inspector, Size to Fit Image makes the object just big enough to
      show the whole image at its original size.
  \end{description}

  \section{The Inspectors Menu}

  The Inspectors menu provides ways for you to interact with the Inspectors sidebar, the floating
  window, or individual palettes and other workspaces you create.

  \begin{description}[style=nextline]
    \item[{Show/Hide Inspectors (\textbf{Shift-Command--I})}] Toggles the Inspectors sidebar to
      the right of the Canvas.
    \item[{Reset to Default}] Choosing this option reverts any changes you have made to the way
      the inspectors are viewed, restoring the inspector sidebar to its default state.
    \item[{Sidebar}] View the inspectors in the sidebar along the right side of OmniGraffle's
      window, as follows:

      \begin{description}[style=nextline]
        \item[{Object (\textbf{Command--1})}] Switches the Inspectors sidebar to reveal the
          following Object inspectors:

          \begin{itemize}
            \item Fill
            \item Stroke
            \item Shadow
            \item Shape
            \item Line
            \item Image
            \item Font
            \item Text Position
            \item Geometry
            \item Alignment
          \end{itemize}
        \item[{Properties (\textbf{Command--2})}] Switches the Inspectors sidebar to reveal the
          following Properties inspectors:

          \begin{itemize}
            \item Connections
            \item Action (\textbf{PRO})
            \item Object Data (\textbf{PRO})
          \end{itemize}
        \item[{Canvas (\textbf{Command--3})}] Switches the Inspectors sidebar to reveal the
          following Canvas inspectors:

          \begin{itemize}
            \item Canvas Size
            \item Canvas Fill
            \item Background Image
            \item Units
            \item Grid
            \item Diagram Layout and Style
            \item Canvas Data
          \end{itemize}
        \item[{Document (\textbf{Command--4})}] Switches the Inspectors sidebar to reveal the
          following Document inspectors:

          \begin{itemize}
            \item Format
            \item Margins
            \item Document Data
          \end{itemize}
      \end{description}
    \item[{Workspaces}] Workspaces allow you to have more control over how you view and use the
      Inspectors in OmniGraffle. There are three initial Workspaces:

      \begin{description}[style=nextline]
        \item[{Sidebar (\textbf{F2})}] This workspace returns the inspectors back to their original
          state as a sidebar attached to the right side of OmniGraffle's window.
        \item[{Floating (\textbf{F3})}] This workspace detaches OmniGraffle's inspectors from the
          main window, creating a floating palette that you can move around on your display---or to
          another display.
        \item[{Palettes (\textbf{F4})}] When selected, this option creates five individual floating
          palettes; one for \textit{each} inspector---Object, Properties, Canvas, Document, and
          Stencils. When the windows are created, they appear in a stack on top of one another. Just
          move them apart, or close the ones you don't need, if that's the case.
      \end{description}
    \item[{}] However, you can create custom Workspaces by ``tearing off'' inspectors into a new
      floating window.
  \end{description}

  \begin{ognote}
    Since most of the \textit{function keys} (for example, F1, F2, F3, and so on) are already
    assigned by the system to function with your Mac hardware, you will need to press the
    \textbf{fn} key and then press the applicable function key. For example, for the Default
    Workspace, you will need to press \textbf{fn-F2}.
  \end{ognote}

  \begin{description}[style=nextline]
    \item[{Edit Workspaces}] Open the Workspaces window, where you can save your favorite
      arrangements of inspectors. Each workspace you save appears at the bottom of this submenu
      and has a function key assigned to it.
  \end{description}

  \section{The Automation Menu (\texorpdfstring{\textbf{PRO}}{PRO})}

  \begin{ogpro}
    OmniGraffle Pro adds \textbf{Omni Automation}, a set of JavaScript-based automation tools for
    scripting interactions within and between OmniGraffle and other applications. The top part of
    this menu contains a list of any Omni Automation Plug-Ins that you have installed; the menu
    items below appear after this list.

    \begin{ognote}
      The primary documentation resource for Omni Automation is maintained at
      \href{https://www.omni-automation.com}{https://www.omni-automation.com}, which provides
      detailed instructions for
      \href{https://omni-automation.com/tutorial/og-macos/index.html}{scripting OmniGraffle} and
      \href{https://omni-automation.com/plugins/index.html}{creating Plug-Ins} as well as
      \href{https://omni-automation.com/plugins/installation.html}{installing Plug-Ins} created by
      others.
    \end{ognote}

    \begin{description}[style=nextline]
      \item[{Show Console}] Opens the Automation Console.
      \item[{Clear Console (\textbf{Command-K})}] Clears the Automation Console of any previous
        input.
      \item[{Console Text Size}] Choose the size for text in the Automation Console from among
        \textbf{Regular}, \textbf{Medium}, and \textbf{Large}.
      \item[{API Reference}] Learn more about what you can do with Omni Automation in OmniGraffle
        by viewing the available classes along with their constructors, functions, and properties.
      \item[{Configure}] Opens the Automation Configuration window, where you can add, remove, and
        manage your Omni Automation Plug-Ins, and set your security preferences for interaction
        with external applications.
    \end{description}

    To learn more about scripting OmniGraffle, visit \href{https://www.omni-automation.com}{the
    Omni Automation website}, or
    \href{https://discourse.omnigroup.com/c/omnigraffle/omnigraffle-automation}{join the
    discussion on our forums}.
  \end{ogpro}

  \section{The Window Menu}

  Commands for working with OmniGraffle's windows.

  \begin{description}[style=nextline]
    \item[{Minimize (\textbf{Command-M})}] Shrink the frontmost window into the Dock.
    \item[{Zoom}] Size the window to fit its contents, or return to the previous window size.
  \end{description}

  \begin{ogpro}
    \begin{description}[style=nextline]
      \item[{New Window on \textit{filename} (\textbf{PRO})}] Opens an additional window of the
        same document, which can be used to independently view different portions of the canvas
        (or different canvases) of the same document.

        The different windows are denoted by the addition of a number in parenthesis next to the
        file name at the top of the toolbar. The numbers are added so you can tell which window is
        which in the Window menu.

        If you place two windows next to each other (or move one of the windows to another display)
        and move an object on the canvas, you will see the object move in the other window.
    \end{description}
  \end{ogpro}

  \begin{ognote}
    The next four menu options are only available when using OmniGraffle 7 on macOS Sierra (v10.12
    or higher). See Show/Hide Tab Bar for more detail.
  \end{ognote}

  \begin{description}[style=nextline]
    \item[{Show Previous Tab (\textbf{Control-Shift-Tab})}] Switches the view to the document in
      the previous tab.
    \item[{Show Next Tab (\textbf{Control-Tab})}] Switches the view to the document in the next
      tab.
    \item[{Move Tab to New Window}] Moves the document in the current tab to a separate window.
    \item[{Merge All Windows}] Collects all open OmniGraffle windows as tabs in the current
      document window.
    \item[{Stencils}] Opens the Stencil Browser so you can drag stencils to the canvas or search
      your system and Stenciltown for the stencil you need.
    \item[{Bring All to Front}] Put all of OmniGraffle's windows in front of all other app
      windows.
  \end{description}

  \section{The Help Menu}

  When you're in a pinch and can't sort something out, we're here to help.

  \begin{description}[style=nextline]
    \item[{OmniGraffle Help}] This summons the onscreen help you're reading right now. Thanks for
      being the kind of person who reads the documentation.
    \item[{Welcome to OmniGraffle 7}] This opens a \textit{Welcome} document, which highlights the
      new features you'll find in OmniGraffle 7.
    \item[{Release Notes}] This pops open a window that includes lots of information from our
      engineering team about the things they've fixed, updated, or have added to OmniGraffle. This
      document changes whenever we issue an update of OmniGraffle, so be sure to read this after
      updating to the latest and greatest build to see what's changed.
    \item[{Sign Up for the Omni Newsletter}] Choosing this item takes you to Omni's website in
      your default web browser, where you can sign up to receive a quarterly email from Omni.
    \item[{Get More Stencils}] Choosing this menu option takes you out of OmniGraffle and opens
      your favorite web browser to \href{https://stenciltown.omnigroup.com}{Stenciltown}, our
      online repository for sharing and discovering new stencils for OmniGraffle.
    \item[{Open Scripts Folder}] This opens a Finder window, taking you to the \textit{Application
      Scripts} folder, used to contain OmniGraffle-related AppleScripts you have created.
    \item[{Contact Omni}] Choose this to start an email message to the OmniGraffle support folks.
      We also give you the option of attaching the OmniGraffle file you are working on, just in
      case there's a bug or you have a particular question how to do something. This is the best
      way to get in touch with us; your version number and license information are automatically
      entered in the subject field.
  \end{description}
  %%%%%%%%%%%%%%%%%%%%%%%%%%%%%%%%%%%%%%%%%%%%%%%%%%%%%%%%%%%%%%%%%%%%%%%%%%%%%%%%%%%%%%%%%%%%%%%%%%%
  \end{document}
